
\chapter{Übungen, Rätsel und ein Rollenspiel}

\section*{1. Buchstaben-Zuordnung}

Zu welchem Buchstaben (\arbig{ب}, \arbig{ت}, \arbig{ث}, \arbig{ن},
\arbig{ي}, \arbig{ا}, \arbig{و}, \arbig{ر}, \arbig{ز} oder \arbig{د})
gehört der \emph{erste} Laut in diesen Wörtern? Schreib den Buchstaben daneben.

\begin{enumerate}[leftmargin=*]
\item \textarabic{\Large بيت} (bayt, Haus) -- \rule{2cm}{0.4pt}
\item \textarabic{\Large تفاح} (tuffaH, Apfel) -- \rule{2cm}{0.4pt}
\item \textarabic{\Large نور} (nuur, Licht) -- \rule{2cm}{0.4pt}
\item \textarabic{\Large ورد} (ward, Rose) -- \rule{2cm}{0.4pt}
\item \textarabic{\Large موز} (mooz, Banane) -- \rule{2cm}{0.4pt}
\item \textarabic{\Large دار} (daar, Haus/Anwesen) -- \rule{2cm}{0.4pt}
\end{enumerate}

\section*{2. Umschrift-Lücken}

Fülle die fehlende Umschrift aus dem Gedächtnis aus:

\begin{enumerate}[leftmargin=*]
\item \textarabic{\Large شكراً} = \rule{3cm}{0.4pt} (Danke)
\item \textarabic{\Large يلا} = \rule{3cm}{0.4pt} (Los geht's!)
\item \textarabic{\Large خلاص} = \rule{3cm}{0.4pt} (Fertig/Genug)
\item \textarabic{\Large بكام ده؟} = \rule{3cm}{0.4pt} (Wie viel kostet das?)
\item \textarabic{\Large ماشي} = \rule{3cm}{0.4pt} (Okay)
\end{enumerate}

\section*{3. Buchstaben-Bingo}

Kreise beim Üben (oder wenn Mohamed Buchstaben ansagt) die passenden Felder ein:

\begin{center}
\renewcommand{\arraystretch}{2.2}
\begin{tabular}{|c|c|c|c|}
\hline
\arbig{ب} & \arbig{ن} & \arbig{و} & \arbig{ذ} \\
\hline
\arbig{ت} & \arbig{ا} & \arbig{ز} & \arbig{ي} \\
\hline
\arbig{د} & \arbig{ث} & \arbig{ر} & \arbig{ن} \\
\hline
\arbig{ي} & \arbig{و} & \arbig{ب} & \arbig{ا} \\
\hline
\end{tabular}
\end{center}

\section*{4. Rollenspiel: ,,Ertappt in Kairo``}

\begin{ammahmoudbox}
\textbf{Am Mahmoud:} \textarabic{أهلاً! إزيك؟} (Willkommen! Wie geht's?)\\
\textbf{Julia:} \textarabic{كويسة، الحمد لله! إنتي إزيك؟} \\
\textbf{Am Mahmoud:} \textarabic{الحمد لله! فين عايزة تروحي؟} (Wohin willst du?) \\
\textbf{Julia:} \textarabic{للسوق، من فضلك.} (Zum Markt, bitte.) \\
\textbf{Am Mahmoud:} \textarabic{ماشي! يلا بينا!} (Okay! Auf geht's!) \\
\textit{-- Am Ziel angekommen --} \\
\textbf{Am Mahmoud:} \textarabic{بمية جنيه.} (Das macht 100 Pfund.) \\
\textbf{Julia:} \textarabic{غالي أوي! خمسين!} (Viel zu teuer! Fünfzig!) \\
\textbf{Am Mahmoud:} \textarabic{طيب... معلش، ماشي!} (Na gut... macht nichts, abgemacht!)
\end{ammahmoudbox}

\begin{bissobox}
Während die beiden verhandeln, klaut Bisso in aller Ruhe ein Stück
Fladenbrot aus Julias Tasche. Niemand hat es kommen sehen. So ist das Leben
in Kairo: man verhandelt um fünfzig Pfund und verliert das Brot trotzdem.
\end{bissobox}

\section*{Lösungen (nicht schummeln, Julia!)}

\begin{center}
\scalebox{-1}[-1]{%
\begin{minipage}{10cm}
\small
1. ب -- 2. ت -- 3. ن -- 4. و -- 5. م -- 6. د\\
shukran -- yalla -- khalaS -- bikaam da? -- maashi
\end{minipage}%
}
\end{center}

\begin{checkpointbox}
\textbf{Geschafft!} Du kennst jetzt elf Buchstaben, alle Vokalzeichen,
eine komplette Begrüßungs-Choreografie, wie man in Kairo feilscht, die
IBM-Regel und die Zahlen 1 bis 10. Das ist ehrlich schon eine Menge --
\textarabic{مبروك!} (\textit{mabruk} -- Glückwunsch!)

Band 2 nimmt die restlichen Buchstaben vor, dazu Café-Bestellungen und
mehr Am Mahmoud. Sag einfach Mohamed Bescheid, wenn du bereit bist.
\end{checkpointbox}
